\documentclass[11pt]{article}
\usepackage[margin=28mm]{geometry}
\usepackage{amsmath,amssymb,amsthm,mathtools,mathrsfs,booktabs,longtable,array}
\usepackage[colorlinks]{hyperref}
\newcommand{\statement}[3]{\subsection*{#1. #3 [#2]}}
\begin{document}
\statement{C.32}{Proved}{The relative pro-$p$ braid monoid has unit parameter}
Fix a prime $p$. Suppose that
$\mu\in1+2\mathbb Z_p$ and $f\in(F_2^{(p)})'$ satisfy the duality,
hexagon and pentagon equations in the corresponding pro-$p$ pure
groups. Then $\mu\in\mathbb Z_p^\times$. In particular the pro-$p$
Grothendieck--Teichm\"uller endomorphism monoid is its unit group.

\begin{proof}
The case $p=2$ follows already from $\mu\in1+2\mathbb Z_2$.
We give the relative braid argument, including its topology and
cohomology, for every prime. Define
\[
 B_n^{[p]}=\widehat B_n/
       \ker(\widehat P_n\longrightarrow P_n^{(p)}).
\]
The kernel being divided out is characteristic in $\widehat P_n$
and therefore normal in $\widehat B_n$. Thus
\[
 1\longrightarrow P_n^{(p)}\longrightarrow B_n^{[p]}
       \longrightarrow S_n\longrightarrow1                    \tag{C.32.1}
\]
is exact. This is a relative completion retaining $S_n$, not the
maximal pro-$p$ quotient of $B_n$.

The Fadell--Neuwirth action of $P_{n-1}$ on its free kernel
$F_{n-1}$ fixes every meridian in abelianization. This is the
almost-direct-product structure recorded explicitly in
\cite[Example 1.1]{CECohen2010}. Given a characteristic finite
$p$-group quotient $F_{n-1}/W$, the image of $P_{n-1}$ in its
automorphism group acts trivially on the Frattini quotient and
is therefore a finite $p$-group. The resulting finite semidirect
product is a finite $p$-group quotient of $P_n$ detecting
$F_{n-1}/W$. Characteristic finite $p$-quotients are cofinal for a
finitely generated free group: intersect the finitely many
automorphic conjugates of a normal subgroup of a given finite
$p$-power index. Consequently the free kernel has its full
induced pro-$p$ topology and
\[
 1\longrightarrow F_{n-1}^{(p)}\longrightarrow P_n^{(p)}
       \longrightarrow P_{n-1}^{(p)}\longrightarrow1             \tag{C.32.2}
\]
is exact and split.

Here ``$p$-good'' will mean that cohomology comparison is an
isomorphism for every finite $p$-primary module whose action extends
continuously to the pro-$p$ completion. Finite-rank free groups have
this property: their degree-zero and degree-one comparisons follow
from the free universal property, and both higher cohomologies
vanish. In (C.32.2), free-kernel cohomology with a finite such module
is finite and vanishes above degree one. Its continuous counterpart
is a finite continuous module for $P_{n-1}^{(p)}$. The
Hochschild--Serre spectral-sequence comparison therefore proves
inductively that every $P_n$ is $p$-good in precisely this sense.

Let $S$ be a Sylow $p$-subgroup of $S_n$, and let $G$ be its
preimage in the discrete group $B_n$. The normal subgroup $P_n$
has $p$-power index in $G$. If $W\triangleleft P_n$ has finite
$p$-power index, its finitely many $G/P_n$ conjugates have an
intersection normal in $G$, still of $p$-power index. Hence the
pro-$p$ topology of $G$ induces the full pro-$p$ topology on $P_n$,
and
\[
 1\longrightarrow P_n^{(p)}\longrightarrow G^{(p)}
       \longrightarrow S\longrightarrow1                       \tag{C.32.3}
\]
is exact. Comparing spectral sequences with finite quotient $S$
now proves $p$-goodness of $G$. Its configuration-space cover is
an aspherical manifold of real dimension $2n$, so
$\operatorname{cd}(G)\leq2n$. Goodness, followed by the finite-module
filtered-colimit argument of C.31, gives
\[
                    \operatorname{cd}_p(G^{(p)})\leq2n.
\]
In particular $G^{(p)}$ has no element of order $p$. The preimage
of $S$ in (C.32.1) is a pro-$p$ group and is exactly $G^{(p)}$:
the natural map is the identity on the full kernel $P_n^{(p)}$
and on the quotient $S$ in (C.32.3). Thus every Sylow $p$-subgroup
of $B_n^{[p]}$ is torsion-free.

Complete the automorphism groups of the parenthesized-braid
groupoids pro-$p$, retaining their labelled object sets. Operadic
composition extends because pro-$p$ completion commutes with finite
products. The groupoid presentation in
\cite[Lemma 7.1 and Proposition 7.3]{CEBHR2019} gives an actual
object-fixing endomorphism from our pair: the elementary braiding
goes to $\beta(\beta^2)^{(\mu-1)/2}$ and the associativity arrow
has parameter $f$. Its defining relations are exactly the two
hexagons and pentagon; the homomorphisms on the pure groups factor
through their pro-$p$ completions since the target pure groups
are pro-$p$. Taking the free label-orbit quotient now gives
endomorphisms $\phi_n$ of $B_n^{[p]}$ preserving $S_n$.
The same ribbon presentation and cabling computation as in C.31
apply over $\mathbb Z_p$ and give
\[
                             \phi_n(z_n)=z_n^\mu.                \tag{C.32.4}
\]
Only powers of pure braids have $\mathbb Z_p$ exponents; the formula
for the elementary braiding above explicitly retains its odd
permutation class, even when $p$ is odd and $\mu=0$.

If $\mu=p\nu$, take $n=p$ and
$a=\phi_p(\delta_p)z_p^{-\nu}$, where
$\delta_p=\sigma_1\cdots\sigma_{p-1}$ and $\delta_p^p=z_p$.
Equation (C.32.4) gives $a^p=1$, while its permutation is a
nontrivial $p$-cycle. Choosing $S$ to contain that cycle places
$a$ in the torsion-free group $G^{(p)}$, a contradiction.
Thus $\mu$ is a unit. The free pro-$p$ Frattini criterion makes
the normalized free-group endomorphism an automorphism, and
closedness of the solution monoid gives its inverse within the
same solution set, as in C.30.
\end{proof}

\statement{C.33}{Proved}{An integral pro-$p$ pentagon has a square-unit parameter}
Let $f\in F_2^{(p)}$ satisfy the pentagon in $P_4^{(p)}$, and use
the convention $[x,y]=xyx^{-1}y^{-1}$ for its coefficient $c_2(f)$.
Then
\[
               D=1+24c_2(f)\in(\mathbb Z_p^\times)^2.            \tag{C.33.1}
\]
For either root $\mu\in\mathbb Z_p^\times$, the pair $(\mu,f)$
satisfies the actual pro-$p$ hexagons and belongs to the pro-$p$
Grothendieck--Teichm\"uller group. In particular the canonical
full pentagon word in A.37 has a peripheral conjugator with that
prescribed word at every maximal pro-$p$ quotient.

\begin{proof}
Furusho's theorem makes $f$ a commutator word and supplies duality
and the hexagon in the free prounipotent group over
$E=\mathbb Q_p(\mu)$, where $\mu^2=D$
\cite[Theorem 2 and Corollary 12]{Furusho2010}. For $p=2$,
$D\equiv1\pmod8$, and for $p=3$, $D\equiv1\pmod3$; these are
already square units. Suppose henceforth that $p\geq5$ and
$D\ne0$. Both $\mu$ and $m=(\mu-1)/2$ belong to the valuation
ring $\mathcal O_E$. We shall prove that $D$ is a nonzero square
modulo $p$ using integral metabelian coefficients. No integrality
of logarithmic or gamma-series coefficients is assumed.

Put $\Lambda=\mathbb Z_p[[A,B]]$, $a=1+A$, $b=1+B$. In the
pro-$p$ semidirect product of the additive module
$\Lambda e_1\oplus\Lambda e_2$ by $\mathbb Z_p^2$, write
\[
 (d,v)(e,w)=(de,v+dw),\qquad
 x=(a,e_1),\quad y=(b,e_2).
\]
The diagonal elements are $d=a^s b^t$ with $s,t\in\mathbb Z_p$.
The commutator is
\[
 [x,y]=(1,w_0),\qquad
 w_0=(1-b)e_1+(a-1)e_2.
\]
There is a unique series $F(a,b)\in\Lambda$ such that the image
of $f$ is $(1,F(a,b)w_0)$. Indeed the closed commutator subgroup
is normally generated by $[x,y]$, conjugation multiplies its
translation by $a^s b^t$, and the module $\Lambda w_0$ is compact
and therefore closed. Uniqueness follows because $\Lambda$ is a
domain and $w_0$ has a nonzero coordinate. This constructs the
needed coefficient; no assertion identifying arbitrary closed
subgroups or the whole derived module is needed. Duality gives
\[
                              F(a,b)=F(b,a).                    \tag{C.33.2}
\]
For arbitrary affine matrices $(d,v)$ and $(e,w)$ the evaluated
commutator has translation $(1-e)v+(d-1)w$, and the evaluated
word $f$ has this translation multiplied by $F(d,e)$. The formula
holds first for discrete words and then by continuity. Its image
under the canonical embedding in the prounipotent group is the
same Magnus expansion, coefficient by coefficient.

Set now $a=1+T$ and evaluate the prounipotent identities at
\[
 x=(a,T),\quad y=(a^{-1},0),\quad
 xy=(1,T),\quad z=(xy)^{-1}=(1,-T).                             \tag{C.33.3}
\]
These matrices lie in $1+T M_2(E[[T]])$, so logarithms and every
prounipotent substitution are defined. Write
\[
 t(a)=(1-a^{-1})F(a,a^{-1}),\qquad V(a)=(a-1)F(1,a).
\]
The three word evaluations and the powers in the hexagon are
\[
 \begin{array}{lll}
 f(x,y)=(1,Tt(a)),& f(z,x)=(1,TV(a)),&
 f(y,z)=(1,-TV(a^{-1})),\\[2pt]
 x^m=\bigl(a^m,T(a^m-1)/(a-1)\bigr),&
 z^m=(1,-mT),&y^m=(a^{-m},0).
 \end{array}                                                    \tag{C.33.4}
\]
The last word evaluation uses (C.33.2). Multiply the six factors
in the hexagon
\[
                    f(x,y)x^m f(z,x)z^m f(y,z)y^m=1
\]
and divide the resulting translation by $Ta^m$. This gives
\[
 a^{-m}t(a)+\frac{1-a^{-m}}{a-1}
                  +V(a)-V(a^{-1})-m=0.                         \tag{C.33.5}
\]
Symmetry gives $t(a^{-1})=-a t(a)$. Substituting $a^{-1}$ for
$a$ in (C.33.5) therefore gives
\[
 -a^{m+1}t(a)+\frac{a(a^m-1)}{a-1}
                  +V(a^{-1})-V(a)-m=0.                         \tag{C.33.6}
\]
Add these identities and put $d=a^{-m}-a^{m+1}$. Since
$1-a^{-m}+a^{m+1}-a=1-a-d$, their sum is
\[
                         d\bigl(t(a)-T^{-1}\bigr)=\mu.
\]
Here $d\ne0$ because $\mu\ne0$, so division in $E((T))$ is
legitimate. It yields the exact integral identity
\[
 \begin{split}
 R(T)&:=1-T(1-a^{-1})F(a,a^{-1})\\
     &=\frac{\mu a^{(\mu-1)/2}T}{a^\mu-1}
             \ \in\ 1+T^2\mathbb Z_p[[T]].
 \end{split}                                                    \tag{C.33.7}
\]

Let $H_\mu(T)=(a^\mu-1)/(\mu T)$, whose constant term is one.
Equation (C.33.7) says $R(T)H_\mu(T)=a^m$. For
$0\leq j\leq p-2$ its coefficients
\[
 [T^j]H_\mu(T)=\frac{(\mu-1)\cdots(\mu-j)}{(j+1)!}
\]
are in $\mathcal O_E$. Also
$[T^{p-1}]a^m=\binom{m}{p-1}$ is integral, since $(p-1)!$ is
a $p$-adic unit. Comparing the coefficient of $T^{p-1}$ in the
product, using $R(T)\in\mathbb Z_p[[T]]$, forces
\[
     [T^{p-1}]H_\mu(T)
       =\frac{\prod_{j=1}^{p-1}(\mu-j)}{p!}\in\mathcal O_E.
\]
For $a^m$ only degrees at most $p-1$ were used, whose factorial
denominators are prime to $p$. Integrality of that series in every
degree was not assumed. Thus
\[
 \mu^{p-1}-1=D^{(p-1)/2}-1\in p\mathcal O_E\cap\mathbb Q_p
                                      =p\mathbb Z_p,
\]
because $\prod_{j=1}^{p-1}(X-j)\equiv X^{p-1}-1\pmod p$.
Euler's criterion and Hensel's lemma now prove (C.33.1).

It remains to exclude $D=0$, rather than taking a limit of
(C.33.7) with the word $f$ held fixed. In that case Furusho's
root is the actual integral parameter $\mu=0$. Since $p\geq5$,
also $m=-1/2\in\mathbb Z_p$. The prounipotent hexagons therefore
descend to the free pro-$p$ group under its injective canonical
embedding; the given pentagon remains a pro-$p$ identity.
C.32 then excludes this nonunit parameter. The same injective
comparison for the square-unit roots already obtained shows that
the hexagons hold in the actual pro-$p$ group. Frattini and the
closed solution-monoid argument give the final unit assertion.
For the canonical word of A.37, its intertwining identity is
already a full identity and hence descends to every pro-$p$
quotient, giving the stated conjugators.
\end{proof}

\paragraph{Literature comparison and remaining full equation.}
Furusho explicitly allowed a quadratic extension for the root and
observed that it might fail to be a $p$-adic unit
\cite[paragraph following Corollary 13]{Furusho2010}.
C.32--C.33 address the integral pro-$p$ word using relative braid
torsion and the displayed Magnus coefficient. The exact full and
relative completions and all coefficient comparisons have been
given, and no literature priority is asserted. These local roots
assemble to a full unit in $\widehat{\mathbb Z}$, but the resulting
hexagon is presently established only in every maximal pro-$p$
quotient. A.39's mixed-prime example explains why this alone does
not imply the full profinite third-cusp equation or marked
involution conjugacy.

\end{document}
